% STS Independent Project - IT systems 2023
% Author(s):	Xin Shen
% 				Department of Information Technology, 
%               Division of Computing Science
% 
% Created:		2023/03/30  (Xin Shen)
%---------------------------------------------------------------------------
\documentclass[12pt,twoside,a4paper,fleqn]{report}
\usepackage{STSStyle}
%---------------------------------------------------------------------------
% Page header (Please don't change)
\setlength{\parindent}{0em}                 % Disables paragraph indentation
\rhead[\nouppercase{\rightmark}]{\thepage}  % Special headings
\lhead[\thepage]{\nouppercase{\leftmark}}   % Special headings
\cfoot{}
%---------------------------------------------------------------------------
% Title page (please fill in)   
\begin{document}
\pagestyle{empty}
\begin{center}
    \parbox[c][\textheight][t]{\textwidth}{
        \vspace{-2cm}
        \begin{center}
            \includesvg[width=4cm]{./figures/UU logo.svg}\\ % Please don't change the university logo
            \vspace{1cm}
            {\Huge \bf Thesis Title}\\  % The main title of the thesis goes here (the main title could be your app's name)
            \vspace{0.5cm}
            \includegraphics[width=5cm]{./figures/App logo.png}\\% Your logo goes here
            \vspace{0.5cm}
            {\Large Thesis Subtitle}\\  % Write a subtitle for the thesis, for example, An app for teamwork
            \vspace{0.5cm}
            {Independent Project in Sociotechnical Systems Engineering -- IT Systems}\\% Please don't change this
            \vspace{1cm}
            \begin{minipage}[t]{0.49\textwidth}
                \centering 
                {\bf \large{Team members:}}\\   % Add your name here, names can be sorted alphabetically by surname
                {Name Surname 1}\\
                {Name Surname 2}\\
                {Name Surname 3}\\
                {Name Surname 4}\\
                {Name Surname 5}
            \end{minipage}
            \begin{minipage}[t]{0.49\textwidth}
                \centering
                {\bf \large{Supervisors:}}\\
                {Kamran Forghani}\\ % Please don't change this
                {Second Supervisor Name}    % Add your second supervisor (TA) name here
            \end{minipage}
            \vfill  
            \today
        \end{center}
    }
\end{center}
\cleardoublepage
%---------------------------------------------------------------------------
% Preamble
\pagenumbering{roman} 				% Begin Roman page numbering (i,ii,...)
%---------------------------------------------------------------------------
% Abstract
\chapter*{Abstract}
 \addcontentsline{toc}{chapter}{Abstract}
    %The abstract summarizes in 300--400 words the entire thesis and follows the thesis structure. Start with what motivates the project and why the rationale for the project is important. State clearly what the goal of the project is and unpack the goal into 3-5 outcomes or objectives of the development process. This should represent half of the abstract. The remaining half briefly describes the approach taken to develop the system and the results obtained in the end. The abstract summarizes the evaluation findings and concludes with the contribution of the project.

    The thesis reports on the work you have done over the term. It will include a description of your project, its design, its users, the planning, and so forth. As a rough guideline, each student will contribute equally to the main content of the thesis, not including the Appendix describing in more detail individual contributions.

 \cleardoublepage
%---------------------------------------------------------------------------
% Acknowledgements
% You can REMOVE this chapter if you don't have any person to acknowledge
\chapter*{Acknowledgements}
\addcontentsline{toc}{chapter}{Acknowledgements}
    The acknowledgments section is your opportunity to thank those who have helped and supported you personally and professionally during your thesis or dissertation process. In your acknowledgments, it's okay to use a more informal style than is usually permitted in academic writing, as well as first-person pronouns. Acknowledgments are not considered part of the academic work itself, but rather your chance to write something more personal. 
\cleardoublepage
%---------------------------------------------------------------------------
%The table of contents is automatically generated by the template
\setcounter{tocdepth}{4}            % Sets the number of section levels in the table of contents to 4
\tableofcontents                    % Creates the table of contents
\cleardoublepage                    % Ends the current page
%---------------------------------------------------------------------------
% Chapters
\pagestyle{fancy}               	% Fancy headings
\pagenumbering{arabic}				% Begin arabic page numbering (1,2,...)
\setlength{\parindent}{20pt}        % Sets default paragraph indentation to 20 pt 


% Command \input{filename} imports the contents from filename.tex

% Please remove the following line (\input{chapters/chapter 0.tex}) after finalizing your thesis
\input{chapters/chapter 0.tex}      % Imports the contents from chapter 0.tex  file        

\input{chapters/chapter 1.tex}      % Imports the contents from chapter 1.tex file

\input{chapters/chapter 2.tex}      % Imports the contents from chapter 2.tex file

\input{chapters/chapter 3.tex}      % Imports the contents from chapter 3.tex file

%---------------------------------------------------------------------------
% Bibliography
% References must be added manually to the references.bib file

\bibliographystyle{unsrt}           % applies an unsorted citation style for the references
\bibliography{references.bib}       % loads the references from the references.bib file and creates the Bibliography section   
%---------------------------------------------------------------------------
% Appendix

\appendix                           % Tells Latex that the upcoming chapters should be put in the Appendix
\input{chapters/appendix.tex}       % Imports the contents from appendix.txt file 
%---------------------------------------------------------------------------
\end{document}
